\section*{Data Availability}
\label{sec:data}
This is a named submission (confirmed with the course instructor,
2026-11-13), so the replication package below is the real, public
repository, not an anonymized mirror:
\url{https://github.com/PASS-IITG/cs5187-term-project-dauntlesscrux}.

It contains \tool's source (\path{src/tsck/}) and 19 engine unit tests;
both rules as data (\path{specs/outparam_reset.json},
\path{specs/bio_chain_pop.json}) with Protocol~A's checker and regression
tests; the pinned OpenSSL SHAs and scan scripts (\path{run_scan.sh},
\path{run_protocol_a.sh}); and, under \path{src/tsck/results/}, the
unedited logs, JSON results, tool versions, timing rerun, manual triage
for both protocols, the reachability and caller re-checks, the
commit-history search, and the tool-comparison evidence behind
Table~\ref{tab:rq2-comparison}. \texttt{make check} reruns the offline
checks; \texttt{make scan} reruns every OpenSSL scan from a clean checkout.

\section*{Author Contributions}
\textbf{Rutvij Wagh (230101113), Protocol B lead:} root-cause analysis of
CVE-2022-4450; formulation of the output-pointer reset invariant;
development of the Protocol B rule specification and checker
(\texttt{outparam.py}); validation against the vulnerable and fixed
OpenSSL revisions; evaluation of the unrelated KDF fixes and the Protocol
B results; analysis of flagged sites and their false positives.
\textbf{Udayraj Choudhary (230101107), Protocol A lead:} investigation of
CVE-2023-0215 and the BIO-chain use-after-free; AddressSanitizer
reproduction and vulnerable/fixed revision analysis; formulation of the
BIO push/pop invariant; oversight and validation of the Protocol A
specification, checker (\path{protocol_a.py}), and regression tests;
investigation of the whole-tree results and its 11 false positives.
\textbf{Joint:} selection of the project topic and OpenSSL as the
research subject; design of the shared tree-sitter parsing and CFG
architecture; selection of research questions and evaluation
methodology; integration of both typestate protocols; comparison with
general-purpose static-analysis tools; literature review and
interpretation of results; paper preparation, proofreading, debugging,
and final validation; preparation of the final presentation and
demonstration. Both authors reviewed and take responsibility for the
full paper and replication package, including the sections produced
with generative-AI assistance disclosed below.

\section*{Use of Generative AI}
For Protocol~B, the authors designed the rule and produced all OpenSSL
builds and measurements; a generative AI assistant (Claude; model version
not recorded) helped draft prose from those results. For Protocol~A, the
assistant helped formulate and implement the checker from the authors'
CVE-2023-0215 notes, ran it on an independently cloned OpenSSL tree, and
produced the initial triage of its 11 whole-tree findings; the authors
reproduced this run on a separate machine (Section~\ref{sec:setup}). The
assistant also wrote the analyses in
\path{results/reachability_rq1_candidates.md},
\path{results/unclear_sites_recheck.md}, and
\path{results/historical_recurrence_search.md}. In the final revision,
ChatGPT (GPT-6) verified bibliographic metadata and ACM CCS identifiers and
edited LaTeX/BibTeX files, and Claude condensed the text to fit the page
limit without changing any result. No AI tool altered recorded
experimental outputs. All source excerpts, logs, and reasoning are in the
replication package, and the authors take responsibility for every claim.
